\begin{document}
\maketitle

\begin{abstract}
We attempt the complete scaling-limit theorem for the concrete order-four
charged field
\[
 \Psi_{\lambda,N}(f)
 =\varepsilon_\lambda:\!\exp(i\lambda\cdot\phi_N)\!:(f),
 \qquad
 \lambda=(\omega_s,\omega_f),\quad
 \norm{\lambda}^2=2,\quad h_\lambda=1.
\]
The algebraic limit-exchange statement is proved in full for an arbitrary
finite group: equivariant injective inductive limits commute isometrically
with crossed products, coefficient and fixed-point expectations commute
with the limit, and their Pimsner--Popa constants remain uniform.  Assuming
the charged field and local algebras exist, the identification of the
extension with the holomorphic $(E_8)_1$ net is reduced, with every
hypothesis displayed, to the Kawahigashi--Longo--M\"uger index theorem,
the rank-eight even-unimodular lattice classification, the
Dong--Mason central-charge-eight theorem, and the
Carpi--Kawahigashi--Longo--Weiner VOA--net correspondence.

The uniform Green-form problem is closed here for the edge-projected
frozen collar, with the explicit all-volume constant
$K_G=\pi^2/4$.  The proof combines an exact particle--hole count,
the sharp fundamental-domain sine inequality, and analytic perturbation
of the isolated edge doublet.  For the edge-projected boson modes we
also construct exact compatible dyadic Fock isometries and prove the
charged-field graph defect with rate $N^{-2}$, including adjoints.
The remaining hard lattice-to-field step is microscopic local-algebra
exhaustion.  The frozen computations
verify the proposed two-edge remainder bound
$\norm{R_N}_{\HS}\le2.11$ only for
$N=16,24,32,48,64,96$ (maximum $2.105079358$), not for all $N$.
Consequently the master theorem remains conditional on one analytic
gap box, stated quantitatively at the end.  This is a full proof attempt,
not a claim that finite observations establish a type-$\mathrm{III}$
scaling limit.
\end{abstract}

\tableofcontents

\section{Verdict and logical form}

\begin{theorem}[Charged MMST master theorem: exact proved form]
\label{thm:master}
Let the finite collars, scaling maps, neutral current limit and charged
operators be as in Section~\ref{sec:setup}.  Suppose the scaling algebra
is exhausted by the neutral currents and
the four charged grades as specified in \textup{(ALG-EXH)} there.  Then:
\begin{enumerate}[label=\textup{(L\arabic*)}]
\item For every $f\in C^\infty(S^1)$ there are $p<\infty$ and
  $C_f<\infty$, independent of $N$, such that on the finite-energy core
  \[
    \norm{\Psi_{\lambda,N}(f)\xi}
    \le C_f\norm{(1+L_{0,N})^p\xi}.
  \tag{1.1}\label{eq:L1}
  \]
\item After the canonical scaling identifications,
  $\Psi_{\lambda,N}(f)$ converges strongly on the common finite-energy
  core to a nonzero closable field $\Psi_\lambda(f)$.  The generated local
  algebras converge to a local, conformal extension net
  $I\mapsto\B(I)$ of the neutral net $I\mapsto\A(I)$.
\item The finite $\Z_4$ crossed products, the canonical coefficient
  expectations, and the fixed-point expectations commute with the
  inductive/scaling limit.  The limiting Pimsner--Popa bound is still
  $E(x)\ge\frac14x$ for $x\ge0$.
\item The resulting extension is the strongly local lattice net
  $\A_{V_{E_8}}$, equivalently the level-one $E_8$ current net.  It is
  holomorphic, has $\mu=1$, central charge $8$, vacuum character
  $E_4/\eta^8$, and local algebras equal to the hyperfinite
  type-$\mathrm{III}_1$ factor.
\end{enumerate}
\end{theorem}

\begin{proof}
L1 follows from Theorem~\ref{thm:L1reduction}.  L2 follows from
Theorem~\ref{thm:L2reduction}.  L3 is the unconditional
Theorem~\ref{thm:crossedlimit}.  L4 follows from
Theorem~\ref{thm:identification}.  The Green estimate used by L1 is the
all-volume Theorem~\ref{thm:UGF}; the graph estimate used by L2 is the
all-volume dyadic Theorem~\ref{thm:TELa}.  The condition still used by L2
is precisely \textup{(ALG-EXH)} in the final gap box; it is not silently
replaced by a finite-volume fit.
\end{proof}

\begin{remark}[What is and is not proved]
The implication in Theorem~\ref{thm:master} is proved.  Uniform Green
domination and local Sobolev convergence are proved in
Theorem~\ref{thm:UGF}, and edge-mode graph convergence is proved in
Theorem~\ref{thm:TELa}.  The antecedent \textup{(ALG-EXH)} is not proved
for the frozen microscopic collar.  Therefore
this memorandum does not close
\texttt{SEAM.MMST.TYPEIII.CHARGED.01}.  It proves L3, completes the
classical reduction of L4, closes L1 and the edge-projected TEL-a part
of L2, and isolates the remaining microscopic local-algebra statement.
\end{remark}

\section{Setup: the frozen collar and charged generator}
\label{sec:setup}

\subsection{The finite CAR collar}

For circumference $N$ and fixed open width $N_y=8$, the one-particle
Qi--Wu--Zhang collar is
\[
 h(k_x,k_y)
 =\sin k_x\,\sigma_x+\sin k_y\,\sigma_y+
   (1-\cos k_x-\cos k_y)\sigma_z .
\tag{2.1}\label{eq:qwz}
\]
The $r$th deck sector, $r\in\Z_4$, has boundary phase $i^r$, hence
longitudinal momenta
\[
 p_n^{(r)}=\frac{2\pi}{N}\left(n+\frac r4\right).
\tag{2.2}\label{eq:momentum}
\]
Sixteen Majorana copies produce the carrier central charge
$16\cdot\frac12=8$.  We write $\mathfrak h_N$ for the one-particle space,
$H_N^{(r)}$ for the corresponding one-particle matrix, and
\[
 C_N^{(r)}
 =\chi_{(-\infty,0)}(H_N^{(r)})
  \frac12\chi_{\{0\}}(H_N^{(r)})
\tag{2.3}\label{eq:covariance}
\]
for the particle--hole symmetric covariance.  The finite observable
algebra is a matrix CAR algebra.  Its low-energy GNS space is denoted
$\Hh_N$, and $L_{0,N}$ denotes the shifted, nonnegative edge-energy
operator in the scaling normalization.

\begin{remark}[Necessary correction: there is no uniform full gap]
\label{rem:gapcorrection}
The topological collar has chiral edge levels with spacing $O(N^{-1})$.
Thus a claim of a spectral gap
$\inf_N\operatorname{gap}(H_N)>0$ is false.  The usable hypothesis is a
\emph{uniform bulk gap after removing the edge band}.  If $Q_N$ is the
edge-band projection, the correct statement is
\[
 \operatorname{spec}\bigl((1-Q_N)H_N(1-Q_N)\bigr)
 \cap(-\delta_{\rm bulk},\delta_{\rm bulk})=\varnothing
\tag{2.4}\label{eq:bulkgap}
\]
for one $\delta_{\rm bulk}>0$ independent of $N$.  This controls bulk
resolvents but does not control the massless longitudinal Green function
by itself.
\end{remark}

\subsection{The exact order-four vector}

In rational coordinates,
\[
 \omega_s=\left(\frac12,\frac12,\frac12,\frac12,\frac12\right),
 \qquad
 \omega_f=\left(\frac34,-\frac14,-\frac14,-\frac14\right),
 \qquad
 \lambda=(\omega_s,\omega_f).
\tag{2.5}\label{eq:lambda}
\]
The exact finite calculation gives
\[
 \norm{\omega_s}^2=\frac54,\quad
 \norm{\omega_f}^2=\frac34,\quad
 \norm{\lambda}^2=2,\quad h_\lambda=1,
\tag{2.6}\label{eq:lambda-norm}
\]
and
\[
 [\lambda]^4=[0],\qquad [\lambda]^2=[v].
\tag{2.7}\label{eq:lambda-order}
\]
The norm-two vectors in
$D_5\oplus A_3+k\lambda$, $k=0,1,2,3$, have census
\[
 [52,64,60,64],
\tag{2.8}\label{eq:census}
\]
so the four cosets contain $240$ roots and the two odd grades contain the
full $128$-dimensional spinor contribution.

\subsection{The finite charged field}

Let $\phi_N=(\phi_N^1,\ldots,\phi_N^8)$ be the bosonic current obtained
from the sixteen Majorana copies by edge-projected collar bosonization,
with zero mode fixed by the $\Z_4$ charge sector.  Each complex-fermion
pair gives one boson.  The longitudinal derivative uses the frozen
lattice symbol
\[
 \omega_N(m)=\frac{N}{\pi}\left|\sin\frac{\pi m}{N}\right|,
 \qquad 0<|m|\le\frac N2.
\tag{2.9}\label{eq:latticesymbol}
\]
The edge projection is the analytic top (or bottom) eigenbranch through
the zero-energy crossing; massive bulk-complement modes are not included
in $\phi_N$.  For a fixed normalized
two-cocycle $\varepsilon_\lambda$, define on the finite-energy vectors
\[
 \Psi_{\lambda,N}(f)
 =\sum_{m}\widehat f_m\,
   \varepsilon_\lambda
   :\!\exp(i\lambda\cdot\phi_N)\!:_m .
 \tag{2.10}\label{eq:chargedfield}
\]
Normal ordering is with respect to $C_N^{(0)}$.  Formula
\eqref{eq:chargedfield} is meaningful at every finite $N$.  The theorem
does not follow merely from that fact: the normal-ordering constant and
the graph norm of the exponential can grow with $N$.

\begin{definition}[Scaling convergence used below]
Let $V_N:\Hh_N\to\Hh_\infty$ be the canonical isometries of the chosen
operator-algebraic renormalization system.  A sequence of closable
operators $X_N$ converges on the finite-energy core
\[
 \Dfin=\bigcup_{E<\infty}\id_{[0,E]}(L_0)\Hh_\infty
 \tag{2.11}\label{eq:core}
\]
if for every $\xi\in\Dfin$, represented at sufficiently large scale,
\[
 V_NX_NV_N^*\xi\longrightarrow X\xi,\qquad
 V_NX_N^*V_N^*\xi\longrightarrow X^*\xi.
 \tag{2.12}\label{eq:graphconv}
\]
The adjoint clause is included because bare pointwise convergence of
unbounded fields is too weak to determine local von Neumann algebras.
\end{definition}

\section{Classical inputs, quoted in the forms used}
\label{sec:classical}

\begin{theorem}[Osborne--Stottmeister free-fermion scaling theorem,
form used here]\label{thm:OS}
For the massless chiral lattice fermion systems constructed by
operator-algebraic renormalization, with the regular scaling functions
and compatible ground states specified in their construction, the
renormalized Koo--Saleur approximants converge strongly on the
finite-particle core to the corresponding Virasoro generators.
The smeared approximants converge on that core, and the same argument
applies to level-one Wess--Zumino--Witten currents which are fermion
bilinears.  Their dynamical correlation functions converge uniformly on
compact time intervals under the stated regularity assumptions on the
smearing functions.
\end{theorem}

This is the content used from Osborne--Stottmeister,
\emph{Conformal field theory from lattice fermions},
Theorems 4.11, 4.16, 5.1, 6.1 and 6.3
\cite{OsborneStottmeister}.  It supplies the neutral bilinear current net
and the common stress tensor.  It does not contain
\eqref{eq:chargedfield}: a spin/lattice exponential changes sectors and
is not a fermion bilinear.

\begin{theorem}[Energy bounds and smearing for unitary VOAs]
\label{thm:CKLWenergy}
Let $V$ be a unitary VOA with conformal Hamiltonian $L_0$, and let
$a\in V$ be homogeneous.  The field $Y(a,z)=\sum_na_nz^{-n-d_a}$
is polynomially energy bounded if there exist $M>0$ and integers
$s,p\ge0$ such that
\[
 \norm{a_n b}
 \le M(1+|n|)^s\norm{(1+L_0)^p b}
\tag{3.1}\label{eq:modebound}
\]
for every finite-energy vector $b$.  Under this hypothesis, for every
$f\in C^\infty(S^1)$ the series
$Y(a,f)=\sum_n\widehat f_n a_n$ converges on finite-energy vectors,
defines a closable operator, preserves
$\Hh^\infty=\bigcap_{q\ge0}\operatorname{Dom}(L_0^q)$, and obeys
\[
 \norm{Y(a,f)b}
 \le M\,\norm{f}_{H^{s+1}}\,
       \norm{(1+L_0)^p b}.
\tag{3.2}\label{eq:smearedbound}
\]
If a unitary VOA is generated by homogeneous energy-bounded fields, all
of its fields are energy bounded.  A simple unitary VOA generated by
weight-one affine currents and conformal vectors is strongly local; its
smeared fields generate an irreducible conformal net.
\end{theorem}

This combines CKLW Proposition 6.1, Theorems 6.8, 8.1 and 8.3
\cite{CKLW}, in exactly the direction used below.  The newer extension
theorem also states that every positive-definite even lattice VOA is
energy bounded \cite[Proposition 4.2]{CarpiTomassini}.  These are
\emph{continuum} statements; they do not assert a bound uniform in a
lattice cutoff.

\begin{theorem}[Kawahigashi--Longo--M\"uger]
\label{thm:KLM}
Let $\mathcal N$ be a local irreducible conformal net satisfying the split
property and strong additivity, and suppose its two-interval inclusion
\[
 \mathcal N(E)\subset\widehat{\mathcal N}(E)
 :=\mathcal N(E')'
\]
has finite index.  Then $\mathcal N$ is completely rational, its DHR
category is a nondegenerate finite braided tensor category, and
\[
 \mu_{\mathcal N}
 =[\widehat{\mathcal N}(E):\mathcal N(E)]
 =\sum_{\rho\in\operatorname{Irr}(\mathcal N)}d(\rho)^2.
\tag{3.3}\label{eq:KLMsum}
\]
If $\mathcal N\subset\mathcal M$ is an irreducible finite-index local
extension of completely rational nets, then
\[
 \mu_{\mathcal M}
 =\frac{\mu_{\mathcal N}}{[\mathcal M:\mathcal N]^2}.
\tag{3.4}\label{eq:KLMextension}
\]
In particular $\mu_{\mathcal M}=1$ exactly when the vacuum is the only
irreducible DHR sector.
\end{theorem}

This is the multi-interval/global-index and extension theorem of
Kawahigashi--Longo--M\"uger \cite{KLM}.

\begin{theorem}[Q-systems and local extensions]
\label{thm:Qsystems}
Let $N\subset M$ be a finite-index inclusion of factors.  A conjugate of
the inclusion determines a standard Q-system $(\theta,w,x)$ in
$\operatorname{End}(N)$, and every standard Q-system reconstructs the
inclusion up to isomorphism fixing $N$.  If the endomorphisms are DHR
endomorphisms of a conformal net, the reconstructed extension is local
if and only if the Q-system is commutative with respect to the DHR
braiding.  For a finite outer group action $G\curvearrowright M$,
$N=M^G$ has index $|G|$, group-average expectation, and the regular
Q-system.
\end{theorem}

This is the Longo--Rehren reconstruction in the local-extension form
\cite{LongoRehren}.

\begin{theorem}[Trace class implies split]\label{thm:split}
Let $\mathcal N$ be a conformal net with conformal Hamiltonian $L_0$.
If
\[
 \Tr(e^{-\beta L_0})<\infty\qquad\text{for every }\beta>0,
\tag{3.5}\label{eq:traceclass}
\]
then $\mathcal N$ has the split property.  More quantitatively, the
trace-class condition implies the conformal Buchholz--Wichmann
nuclearity maps are nuclear; the Buchholz--Wichmann condition implies
statistical independence for strictly separated intervals.
\end{theorem}

The implication is due to Buchholz--Wichmann and its conformal
refinements by Buchholz--D'Antoni--Longo \cite{BuchholzWichmann,BDL}.

\begin{theorem}[Central charge eight and the $E_8$ theory]
\label{thm:DM}
Let $V$ be a simple, rational, $C_2$-cofinite, holomorphic unitary VOA of
CFT type and central charge $8$.  Then its character is
\[
 \operatorname{ch}_V(\tau)
 =\frac{E_4(\tau)}{\eta(\tau)^8}
 =q^{-1/3}(1+248q+4124q^2+34752q^3+\cdots),
\tag{3.6}\label{eq:E8char}
\]
$V_1$ is the simple Lie algebra $E_8$ at level one, and
$V\cong V_{E_8}$.  The strongly local net associated by the CKLW
construction is the level-one $E_8$ current net.
\end{theorem}

The VOA classification statement is the central-charge-eight theorem of
Dong--Mason \cite{DongMason}; the last sentence is CKLW.  We stress that
KLM alone proves holomorphy, not the name ``$E_8$''.  The VOA
reconstruction/weight-one argument is the additional identification step.

